% Part: incompleteness
% Chapter: incompleteness-provability
% Section: first-incompleteness-thm

\documentclass[../../../include/open-logic-section]{subfiles}

\begin{document}

\olfileid[hi]{inc}{inp}{1in}

\olsection{प्रथम अपूर्णता-प्रमेय}

अब हम प्रथम अपूर्णता-प्रमेय के ग्योडेल (G\"odel) के मूल प्रमाण का वर्णन
कर सकते हैं। $\Th{T}$ ऐसी भाषा में कोई संगणनीयतः अभिगृहीतीकृत सिद्धांत
हो जो अंकगणित की भाषा का विस्तार करती है, और मान लें कि $\Th{T}$ में
$\Th{Q}$ के अभिगृहीत सम्मिलित हैं। इसका विशेषतः अर्थ है कि $\Th{T}$
संगणनीय फलनों और संबंधों का निरूपण करता है।

हमने तर्क दिया है कि सूत्रों और प्रमाणों का संख्याओं द्वारा उचित कूटलेखन
दिए जाने पर संबंध $\Prf[\Th{T}](x,y)$ संगणनीय है, जहाँ
$\Prf[\Th{T}](x,y)$ तभी और केवल तभी लागू होता है जब $x$,~$\Th{T}$ में
ग्योडेल (G\"odel) संख्या~$y$ वाले !!{formula} की !!a{derivation} की ग्योडेल (G\"odel) संख्या
हो। वास्तव में ग्योडेल (G\"odel) के मन में जो विशेष सिद्धांत था, उसके लिए अपने
शोधपत्र के 45 फलनों और संबंधों की सूची का उपयोग करके ग्योडेल (G\"odel) दिखा सके कि यह
संबंध आदिम पुनरावर्ती है। 45वाँ संबंध $x B y$, उनके~$\Th{T}$ के विशेष
चुनाव के लिए केवल $\Prf[\Th{T}](x,y)$ है। याद रखें कि ग्योडेल (G\"odel) अपने शोधपत्र
में जहाँ ``पुनरावर्ती'' शब्द का उपयोग करते हैं, वहाँ अब हम ``आदिम
पुनरावर्ती'' पदबंध का उपयोग करेंगे।

चूँकि $\Prf[\Th{T}](x,y)$ संगणनीय है, उसका $\Th{T}$ में निरूपण होता है।
उसका निरूपण करने वाले सूत्र को हम $\OPrf[\Th{T}](x,y)$ लिखेंगे।
$\OProv[\Th{T}](y)$ को सूत्र
$\lexists[x][\OPrf[\Th{T}](x,y)]$ मानें। यह ग्योडेल (G\"odel) की सूची के 46वें
संबंध $\fn{Bew}(y)$ का वर्णन करता है। ग्योडेल (G\"odel) के शब्दों में यही एकमात्र
संबंध है जिसके बारे में ``यह नहीं कहा जा सकता कि वह पुनरावर्ती है।'' संभवतः
उनका अर्थ यह था: परिभाषा से यह स्पष्ट नहीं है कि वह संगणनीय है; और बाद के
विकास वास्तव में दिखाते हैं कि वह संगणनीय नहीं है।

$\Th{T}$ ऐसा कोई !!{axiomatizable} सिद्धांत हो जिसमें~$\Th{Q}$ सम्मिलित
है। तब $\Prf[\Th{T}](x, y)$ निर्णेय है, इसलिए उसका~$\Th{Q}$ में किसी
!!a{formula}~$\OPrf[\Th{T}](x, y)$ द्वारा निरूपण होता है।
$\OProv[\Th{T}](y)$ ऊपर वर्णित सूत्र हो। नियत-बिंदु उपपत्ति से कोई सूत्र
$!G_\Th{T}$ ऐसा है कि $\Th{Q}$ (और इसलिए $\Th{T}$) निम्न को !!{derive}
करता है:
\begin{equation}
\ollabel{eqn:qpf}
!G_\Th{T} \liff \lnot \OProv[\Th{T}](\gn{!G_\Th{T}}).
\end{equation}
ध्यान दें कि सारतः $!G_\Th{T}$ कहता है, ``$!G_\Th{T}$,~$\Th{T}$ में
!!{derivable} नहीं है।''

\begin{lem}\ollabel{lem:cons-G-unprov}
यदि $\Th{T}$,~$\Th{Q}$ का कोई संगत, !!{axiomatizable} विस्तार है, तो
$\Th{T} \Proves/ !G_\Th{T}$।
\end{lem}

\begin{proof}
मान लें $\Th{T}$,~$!G_\Th{T}$ को !!{derive} करता है। तब कोई
!!a{derivation} \emph{सचमुच} है, और इसलिए किसी संख्या $m$ के लिए संबंध
$\Prf[\Th{T}](m, \Gn{!G_\Th{T}})$ लागू होता है। किंतु तब $\Th{Q}$ वाक्य
$\OPrf[\Th{T}](\num m, \gn{!G_\Th{T}})$ को !!{derive} करता है। अतः
$\Th{Q}$,~$\lexists[x][\OPrf[\Th{T}](x,\gn{!G_\Th{T}})]$ को !!{derive}
करता है, जो परिभाषा से $\OProv[\Th{T}](\gn{!G_\Th{T}})$ है।
\olref{eqn:qpf} से $\Th{Q}$,~$\lnot !G_\Th{T}$ को !!{derive} करता है;
और चूँकि $\Th{T}$,~$\Th{Q}$ का विस्तार है, $\Th{T}$ भी ऐसा करता है। हमने
दिखाया कि यदि $\Th{T}$,~$!G_\Th{T}$ को !!{derive} करता है, तो वह
$\lnot !G_\Th{T}$ को भी !!{derive} करता है और इसलिए असंगत होगा।
\end{proof}

\begin{defn}
\ollabel{thm:oconsis-q}
कोई सिद्धांत $\Th{T}$ \emph{$\omega$-संगत} है यदि निम्न लागू हो: यदि
$\lexists[x][!A(x)]$ कोई वाक्य है और $\Th{T}$,
$\lnot !A(\num 0)$, $\lnot !A(\num 1)$, $\lnot !A(\num 2)$, \dots को
!!{derive} करता है, तो $\Th{T}$,~$\lexists[x][!A(x)]$ को सिद्ध नहीं करता।
\end{defn}

ध्यान दें कि प्रत्येक $\omega$-संगत सिद्धांत संगत भी है। यह सीधे इस तथ्य
से निकलता है कि यदि $\Th{T}$ असंगत है, तो प्रत्येक~$!A$ के लिए
$\Th{T} \Proves !A$। विशेषतः, यदि $\Th{T}$ असंगत है, तो वह प्रत्येक~$n$
के लिए $\lnot !A(\num n)$ को !!{derive} करता है और
$\lexists[x][!A(x)]$ को भी !!{derive} करता है। अतः यदि $\Th{T}$ असंगत है,
तो वह $\omega$-असंगत है। प्रतिधनात्मक रूप से, यदि $\Th{T}$,
$\omega$-संगत है, तो उसे संगत होना चाहिए।

\begin{lem}\ollabel{lem:omega-cons-G-unref}
यदि $\Th{T}$,~$\Th{Q}$ का कोई $\omega$-संगत, !!{axiomatizable} विस्तार
है, तो $\Th{T} \Proves/ \lnot !G_\Th{T}$।
\end{lem}

\begin{proof}
हम दिखाते हैं कि यदि $\Th{T}$,~$\lnot !G_\Th{T}$ को !!{derive} करता है,
तो वह $\omega$-असंगत है। मान लें $\Th{T}$,~$\lnot !G_\Th{T}$ को
!!{derive} करता है। यदि $\Th{T}$ असंगत है, तो वह $\omega$-असंगत है और
हमारा काम पूरा है। अन्यथा $\Th{T}$ संगत है, इसलिए
\olref{lem:cons-G-unprov} से वह $!G_\Th{T}$ को !!{derive} नहीं करता। चूँकि
$\Th{T}$ में $!G_\Th{T}$ की कोई !!{derivation} नहीं है, $\Th{Q}$ निम्न को
!!{derive} करता है:
\[
\lnot \OPrf[\Th{T}](\num 0, \gn{!G_\Th{T}}), \lnot \OPrf[\Th{T}](\num 1,
\gn{!G_\Th{T}}), \lnot \OPrf[\Th{T}](\num 2, \gn{!G_\Th{T}}), \dots
\]
और $\Th{T}$ भी ऐसा करता है। दूसरी ओर \olref{eqn:qpf} से
$\lnot !G_\Th{T}$,
$\lexists[x][\OPrf[\Th{T}](x,\gn{!G_\Th{T}})]$ के समतुल्य है। अतः
$\Th{T}$,~$\omega$-असंगत है।
\end{proof}

\begin{prob}
  प्रत्येक $\omega$-संगत सिद्धांत संगत है। दिखाएँ कि विलोम लागू नहीं होता,
  अर्थात् ऐसे संगत किंतु $\omega$-असंगत सिद्धांत हैं। यह दिखाकर ऐसा करें
  कि $\Th{Q} \cup \{\lnot !G_\Th{Q}\}$ संगत किंतु $\omega$-असंगत है।
\end{prob}

\begin{thm}
\ollabel{thm:first-incompleteness} $\Th{T}$,~$\Th{Q}$ का कोई
$\omega$-संगत, !!{axiomatizable} विस्तार हो। तब $\Th{T}$ पूर्ण नहीं है।
\end{thm}

\begin{proof}
  यदि $\Th{T}$,~$\omega$-संगत है, तो वह संगत है; अतः
  \olref{lem:cons-G-unprov} से $\Th{T} \Proves/ !G_\Th{T}$।
  \olref{lem:omega-cons-G-unref} से
  $\Th{T} \Proves/ \lnot !G_\Th{T}$। इसका अर्थ है कि $\Th{T}$ अपूर्ण है,
  क्योंकि वह न $!G_\Th{T}$ को और न $\lnot !G_\Th{T}$ को !!{derive} करता
  है।
\end{proof}

\end{document}
